% Fragmento integrable. Preambulo anfitrion: amsmath, amssymb, longtable, hyperref.
\section{Dodecaphase profile and quadratic return}
\label{hmtfg:fragmento:perfil-dodecafasico-y-retorno-cuadratico}

\subsection{Genealogy and object of the proof}
\label{hmtfg:desarrollo:1}

The initial object is the critical readout of the oriented dodecaphase wheel already constructed in the corpus. APP evaluates addition and multiplication on the two sheets of the digit lattice, preserving quotient and remainder; TRIT determines regime and orientation; TPK transport selects, transports and updates the state with its carry, boundary and memory records. The nonadic connection prolongs that state, with phase return and memory advance. Its readers act on the same discrete structure of the continuum.

The focal reader of the source is
\[
\mathcal P_{\mathrm{crit}}:U_{12}^{\mathrm{sign}}\longrightarrow
\mathbb R/2\pi\mathbb Z,\qquad
m\longmapsto\phi_m=\frac{(3m-1)\pi}{18},\quad 1\le m\le12.
\]
To evaluate the angular moment, the oriented representative of that chart is retained: replacing it with another representative modulo \(2\pi\) would change the moment. The uniform trace fixes \(w_m=1/12\). The focal chain is
\[
\text{APP–TRIT–TPK}\longrightarrow
\text{discrete structure of the continuum}\longrightarrow U_{12}^{\mathrm{sign}}
\xrightarrow{\mathcal P_{\mathrm{crit}}}(\phi_m,w_m)
\longrightarrow u_0\longrightarrow f_\lambda
\longrightarrow\text{returns and renormalization}.
\]

The abbreviated notation retains the complete state and its prolongations as antecedent; the scalar family is a readout of that antecedent. The already fixed reader, with its order, orientation and uniform sector, is taken as an explicit starting structure. The conclusions concern that particular reader; uniqueness of that reader among all possible TPK readouts lies outside these proofs.

The coordinate \(\pi_{\mathrm{HMT}}\) used by the chart belongs to the internal outputs of the corpus. Its factor cancels exactly when normalizing the second moment. The Feigenbaum values enter exclusively in subsequent recognition, never in the definition of the family or the certification of roots.


\subsection{Exact profile and quadratic criticality}
\label{hmtfg:desarrollo:2}

Let \(s_m=3m-1\). The integer sum
\[
\sum_{m=1}^{12}s_m^2=5394=6\cdot899
\]
determines
\[
\kappa_0=\frac{36}{\pi\sqrt{899}},\qquad
k_m=\kappa_0\phi_m=\frac{2s_m}{\sqrt{899}},
\]
\[
u_0(x)=\frac1{12}\sum_{m=1}^{12}(1-\cos(k_mx)),\qquad
f_{\lambda,0}(x)=1-\lambda u_0(x).
\]
By the second-moment identity,
\[
u_0(0)=u_0'(0)=0,\quad u_0''(0)=2,\quad
f_{\lambda,0}''(0)=-2\lambda.
\]
The analytic representation preserves rational coefficients:
\[
u_0(z)=\sum_{j\ge1}(-1)^{j+1}
\frac{4^j\sum_m s_m^{2j}}{12\,899^j(2j)!}z^{2j}
=z^2-\frac{715769}{2424603}z^4
+\frac{1353832978}{32695771455}z^6+\cdots.
\]
Thus the finite record of twelve sectors produces an entire profile, with quadratic criticality and scale fixed before iteration.


\subsection{Structural robustness with uniform bounds}
\label{hmtfg:desarrollo:3}

The deformation already present in the originating code is
\[
w_m(\eta)=\frac{1+\eta p_m}{12},\qquad
p_m=12\cos(3\phi_m)-\cos(4\phi_m).
\]
We define
\[
M_{2,\eta}=\sum_mw_m(\eta)\phi_m^2,\quad
\kappa_\eta=\sqrt{2/M_{2,\eta}},\quad k_{m,\eta}=\kappa_\eta\phi_m,
\]
\[
u_\eta(x)=\sum_mw_m(\eta)(1-\cos(k_{m,\eta}x)),\qquad
f_{\lambda,\eta}=1-\lambda u_\eta.
\]

\textbf{Structural robustness theorem.} For \(|\eta|\le1/40\), the weights are positive and sum to one. The profile is even and entire, satisfies \(u_\eta''(0)=2\) and is strictly increasing on \((0,1]\). For \(0<\lambda\le2/u_\eta(1)\), the map sends \([-1,1]\) into itself, has a unique critical point in that interval and satisfies
\[
S f_{\lambda,\eta}(x)\le-\frac{640}{47647}<0,\qquad0<|x|\le1.
\]

\textbf{Proof.} The sums of cosines \(3\phi_m\) and \(4\phi_m\) vanish through cycles of lengths four and three. Since \(|p_m|\le13\),
\[
\frac9{160}\le w_m\le\frac{53}{480},\qquad
\frac{27}{40}M_{2,0}\le M_{2,\eta}\le\frac{53}{40}M_{2,0}.
\]
It follows that
\[
k_{\max,\eta}^2\le\frac{196000}{24273}<9<\pi^2,\qquad
k_{\min,\eta}^2\ge\frac{640}{47647}.
\]
For \(0<x\le1\), all sines \(\sin(k_{m,\eta}x)\) are positive, whence
\[
u_\eta'(x)=\sum_mw_mk_m\sin(k_mx)>0,\qquad
u_\eta'''(x)=-\sum_mw_mk_m^3\sin(k_mx)<0.
\]
The quotient \(u_\eta'''/u_\eta'\) is the negative of a positive mean of the \(k_m^2\). Invariance of the Schwarzian under affine changes of the value gives
\[
S f_{\lambda,\eta}
=\frac{u_\eta'''}{u_\eta'}-\frac32
\left(\frac{u_\eta''}{u_\eta'}\right)^2
\le-k_{\min,\eta}^2.
\]
Evenness covers \(x<0\), and the image of the interval is
\([1-\lambda u_\eta(1),1]\). \(\square\)

The width \(1/40\) is a choice made to obtain uniform bounds, distinct from a fundamental constant or an optimal threshold. The original sector remains \(\eta=0\).


\subsection{Global quadratic holomorphic coordinate on a strip}
\label{hmtfg:desarrollo:4}

\textbf{Holomorphic factorization theorem.} For \(|\eta|\le1/40\), in
\[
\mathcal S=\{z\in\mathbb C:|\operatorname{Re}z|<\pi/3\}
\]
there exists a holomorphic, odd and univalent function \(b_\eta\), with \(b_\eta'(0)=1\), such that
\[
\boxed{u_\eta(z)=b_\eta(z)^2,\qquad
f_{\lambda,\eta}(z)=1-\lambda b_\eta(z)^2.}
\]

\textbf{Proof.} For \(z=x+iy\) in the right half-strip,
\[
\operatorname{Re}u_\eta'(z)
=\sum_mw_mk_m\sin(k_mx)\cosh(k_my)>0.
\]
The integral of \(u_\eta'\) over the segment between two points, divided by the difference of the endpoints, has positive real part. Convexity of the half-strip proves injectivity there. Evenness gives the corresponding result on the left.

Moreover,
\[
\operatorname{Im}u_\eta(x+iy)
=\sum_mw_m\sin(k_mx)\sinh(k_my)
\]
has the sign of \(y\) for \(x>0\). On the positive real axis \(u_\eta>0\); on the imaginary axis,
\[
u_\eta(iy)=\sum_mw_m(1-\cosh(k_my))
\]
is strictly negative away from zero and decreases strictly with \(|y|\). These separations, injectivity on each half-strip and evenness prove that the fibers of \(u_\eta\) in \(\mathcal S\) are exactly \(\{z,-z\}\), except for the fiber of the origin.

The function \(u_\eta(z)/z^2\), extended with value one at zero, is holomorphic and has no zeros in the simply connected strip. It has a holomorphic logarithm \(L_\eta\) with \(L_\eta(0)=0\), which is even by uniqueness of the normalization. Define
\[
b_\eta(z)=z\exp\!\left(\frac12L_\eta(z)\right).
\]
Then \(b_\eta^2=u_\eta\), \(b_\eta\) is odd and \(b_\eta'(0)=1\). The equality \(b_\eta(z)=b_\eta(w)\) forces \(w=z\) or \(w=-z\). The second case gives equality only at the origin. This proves its univalence. \(\square\)

The twelve phases thus determine an explicit conformal deformation of the quadratic fold. The complete dynamical conjugacy preserves that deformation:
\[
b_\eta\circ f_{\lambda,\eta}\circ b_\eta^{-1}(w)
=b_\eta(1-\lambda w^2),
\]
on their composition domains. The factorization proved and a conjugacy to the entire quadratic family are distinct assertions.


\subsection{Return, rescaling and memory}
\label{hmtfg:desarrollo:5}

Let \(a=-f(1)>0\), \(H_a(x)=-ax\), \(J=H_a([-1,1])\). When \(J\subset[-1,1]\) is admissible for the return, \(f^2(J)\subset J\) and the compositions are defined,
\[
\mathcal Rf=H_a^{-1}f^2H_a,\qquad
\mathcal Rf(x)=-a^{-1}f(f(-ax)).
\]
If in addition \(f(J)\subset(0,1]\), a unique central quadratic maximum is preserved:
\[
(\mathcal Rf)(0)=1,\qquad
(\mathcal Rf)''(0)=-a f'(1)f''(0)<0.
\]
Composition of Schwarzians gives, at regular points,
\[
S(\mathcal Rf)(x)=a^2
\left[(Sf)(f(-ax))f'(-ax)^2+(Sf)(-ax)\right]<0.
\]
Domain inclusions are part of the hypotheses and are checked at every scale.

Let \(v_\eta=(u_\eta|_{[0,1]})^{-1}\). The inverse branches are
\[
\beta_\sigma(y)=\sigma v_\eta((1-y)/\lambda),\quad
\sigma\in\{-1,+1\},
\]
with the critical record \(\beta_0(1)=0\). The step
\[
(x,w)\longmapsto(f(x),w\mathbin{\|}\sigma(x))
\]
is inverted on its image by reading the last sheet and applying \(\beta_\sigma\). The word preserves branch decisions; the previous value is reconstructed through the map equation.

The return groups two transitions. From the endpoint \(x'\) and its sheets \((\sigma_0,\sigma_1)\) one recovers
\[
x=H_a^{-1}\beta_{\sigma_0}
\!\left(\beta_{\sigma_1}(H_ax')\right).
\]
The intermediate points are also reconstructed. The APP–TRIT–TPK records of sheet, carry, route, boundary and memory are concatenated while preserving their identity. The sign of the analytic branch is an additional record: it does not replace those data with a binary label.

If \(C_N\) groups histories of length \(2N\) into pairs, the restrictions satisfy
\[
\pi^{\mathrm{ret}}_{N,M}C_N
=C_M\pi^f_{2N,2M},\qquad M\le N.
\]
Both sides remove the same terminal pairs and preserve the same inverse branches. Compatible transport of finite histories is thus proved.


\subsection{Accumulated scales and derivative of the operator}
\label{hmtfg:desarrollo:6}

Let \(c_n=f^{2^n}(0)\), \(c_0=1\) and \(H_n(x)=c_nx\). As long as the compositions are admissible and \(c_n\ne0\),
\[
\boxed{\mathcal R^nf=H_n^{-1}f^{2^n}H_n},\qquad
(\mathcal R^nf)(1)=\frac{c_{n+1}}{c_n}=-a_n.
\]
Induction follows from \(H_nH_{a_n}=H_{n+1}\). The convention \(a_n>0\) requires sign alternation of \(c_n\). At a superstable root of period \(2^N\), \(c_N=0\): that normalization becomes singular and the preceding levels or a controlled limit must be used.

For \(z=-ax\), \(w=f(z)\) and a perturbation \(h\), the complete derivative is
\[
\boxed{
D\mathcal R_f[h](x)=
-\frac{h(w)+f'(w)h(z)}a+
\frac{h(1)}a
\left[\mathcal Rf(x)-xf'(w)f'(z)\right].
}
\]
It is differentiated using \(\dot a=-h(1)\), \(\dot z=xh(1)\),
\(\dot w=h(z)+xf'(z)h(1)\) and the variation of the denominator. If \(h(0)=0\), this gives \(D\mathcal R_f[h](0)=0\). The initial HMT parameter direction is \(h=-u_\eta\).

This formula includes scale variation; freezing \(a\) would give a different linearized operator.


\subsection{Certified roots and itineraries through period 256}
\label{hmtfg:desarrollo:7}

Let \(F_n(\lambda)=f_{\lambda,0}^{2^n}(0)\). Compute
\[
x_0=p_0=0,\qquad x_{j+1}=1-\lambda u_0(x_j),\qquad
p_{j+1}=-u_0(x_j)-\lambda u_0'(x_j)p_j.
\]
Then \(F_n=x_{2^n}\) and \(F_n'=p_{2^n}\).

The attached program uses rational intervals with outward integer rounding, square roots enclosed by integer square roots, and sine/cosine with an explicit Taylor remainder. The exact bound \(|u_0''|\le2\) allows centered evaluation with a quadratic remainder.

The historical approximations only propose boxes of radius \(10^{-42}\). Opposite signs at their endpoints, separation of \(F_n'\) from zero throughout the box, and exclusion of returns with periods that properly divide \(2^n\) are certified. The intermediate value theorem and monotonicity prove local existence and uniqueness; exclusion of divisors proves the least period.

\begingroup\small
\begin{longtable}{p{.12\linewidth}p{.17\linewidth}p{.62\linewidth}}
\hline
Level & Least period & Box center, rounded \\
\hline
1 & 2 & 1.345913990471832792210949 \\
2 & 4 & 1.763379787846910127290794 \\
3 & 8 & 1.850413796950136464530567 \\
4 & 16 & 1.868979756606798125150574 \\
5 & 32 & 1.872955034435659740004205 \\
6 & 64 & 1.873806367467030119412262 \\
7 & 128 & 1.873988695221365362307169 \\
8 & 256 & 1.874027744154450672897008 \\
\hline\end{longtable}
\endgroup

The first root has the exact expression \(\lambda_1=1/u_0(1)\). The complete boxes and their checks are in the JSON certificate.

The 502 preclosure states of these eight levels, all separated from zero, and their complete sign words have been enclosed. In particular,
\[
\operatorname{sign}f_{\lambda_n,0}^{2^j}(0)=(-1)^j,\qquad0\le j<n.
\]
The finite orientations of the rescalings preceding closure are established. Inclusions of restrictive intervals at every scale constitute an additional check.

For \(\delta_{\mathrm F,n}=(\lambda_{n-1}-\lambda_{n-2})/(\lambda_n-\lambda_{n-1})\), an abbreviated outer decimal interval is
\[
\boxed{
4.66921218915020415455609621588966392804
<\delta_{\mathrm F,8}<
4.66921218915020415455609621588966392863.
}
\]
The width of the complete interval is less than \(5.808\cdot10^{-37}\). It measures the numerical error of the finite ratio \(\delta_{\mathrm F,8}\), distinct from the error in approximating the universal limit.


\subsection{Analytic persistence of the deformed roots}
\label{hmtfg:desarrollo:8}

For each certified level, \(F_n'(\lambda_n)\ne0\). The implicit function theorem produces a local analytic branch \(\lambda_n(\eta)\) with
\[
f_{\lambda_n(\eta),\eta}^{2^n}(0)=0.
\]
By continuity, the proper returns remain separated from zero for \(\eta\) sufficiently small; the least period and itinerary are preserved. The existence of those local neighborhoods is proved. Their common width is not yet identified with the entire structural strip \(1/40\).

With \(M=M_{2,\eta}\), the profile response is
\[
v_\eta(x)=\partial_\eta u_\eta(x)
=\sum_m\frac{p_m}{12}(1-\cos(k_{m,\eta}x))
-\frac{M'}{2M}\,x u_\eta'(x).
\]
The second term preserves normalization of the second moment. If \(q_j=\partial_\eta x_j\) at fixed \(\lambda\),
\[
q_0=0,\qquad q_{j+1}=-\lambda v_\eta(x_j)-\lambda u_\eta'(x_j)q_j,
\qquad
\boxed{\lambda_n'(\eta)=-q_{2^n}/p_{2^n}}.
\]
The weights–scale–response–superstable-parameter chain is thus explicit. This finite nondegeneracy is distinct from asymptotic transversality to the stable leaf of the universal fixed point.


\subsection{Analytic tails and error propagation}
\label{hmtfg:desarrollo:9}

For \(\eta=0\), let
\[
\mu_{2j}=\frac1{12}\sum_mk_m^{2j},\quad
p_N(z)=\sum_{j=1}^N(-1)^{j+1}\frac{\mu_{2j}}{(2j)!}z^{2j},
\quad\Omega=\frac{70}{\sqrt{899}}.
\]
If \(q_N=\Omega^2r^2/[(2N+3)(2N+4)]<1\),
\[
\|u_0-p_N\|_r\le
\frac{\mu_{2N+2}r^{2N+2}}{(2N+2)!(1-q_N)}.
\]
The inequality \(\mu_{2j+2}\le\Omega^2\mu_{2j}\) bounds the tail by a geometric series. For \(0\le s\le2N+2\),
\[
\|(u_0-p_N)^{(s)}\|_r\le
\frac{\mu_{2N+2}r^{2N+2-s}}{(2N+2-s)!}
\left(1-\frac{\Omega^2r^2}{(2N+3-s)(2N+4-s)}\right)^{-1},
\]
when the denominator is positive.

To transport errors through a return, its domains are also controlled: if \(\|f-g\|\le\varepsilon\), \(a_f,a_g\ge a_*>0\), the compositions lie in a common convex domain and \(\|f'\|\le L,\ \|g\|\le M\) there, then
\[
\|\mathcal Rf-\mathcal Rg\|_r
\le\varepsilon\left[\frac{1+L+L^2r}{a_*}+\frac{M}{a_*^2}\right].
\]
The proof adds the outer perturbation, the inner perturbation, the change of argument due to \(|a_f-a_g|\le\varepsilon\) and the variation of the denominator. These bounds prevent reuse of the initial cosine sum as if all renormalized profiles retained that same form.
